\documentclass[12pt,a4paper]{article}
\usepackage[brazil]{babel}
\usepackage[utf8]{inputenc}
\usepackage[T1]{fontenc}
\usepackage{geometry}
\geometry{top=3cm, bottom=2.5cm, left=3cm, right=2.5cm}
\usepackage{graphicx}
\usepackage{booktabs}
\usepackage{array}
\usepackage{tabularx}
\usepackage{hyperref}
\hypersetup{colorlinks=true, linkcolor=blue, citecolor=blue, urlcolor=blue, filecolor=blue}
\usepackage{listings}
\usepackage{xcolor}
\usepackage{float}
\usepackage{caption}
\usepackage{enumitem}
\captionsetup[table]{position=top, font=small}
\captionsetup[figure]{position=bottom, font=small}

\lstset{
  basicstyle=\ttfamily\small,
  breaklines=true,
  frame=single,
  backgroundcolor=\color{gray!10},
  keywordstyle=\color{blue},
  commentstyle=\color{gray},
  stringstyle=\color{red!70!black},
  numbers=left,
  numberstyle=\tiny\color{gray},
  stepnumber=1,
  numbersep=5pt
}

\begin{document}

\title{\textbf{Projeto Final: Sistema de Regressão Espectral com Bagging e Subagging}\\
\small Reengenharia e Implementação em Streamlit/Python}
\author{Wagner Oliveira de Araujo}
\date{IFG --- 2026}
\maketitle
\thispagestyle{empty}

\newpage
\tableofcontents
\newpage

\section{Levantamento de Requisitos}

O levantamento visou catalogar, classificar e priorizar as funcionalidades do sistema legado (Vaadin/Java) para orientar a reengenharia em Python/Streamlit, garantindo rastreabilidade e aderência às necessidades de pesquisadores e alunos de quimiometria.

\subsection{Requisitos Funcionais}
\begin{table}[H]
\centering
\caption{Inventário das funções do sistema}
\begin{tabularx}{\textwidth}{l l X}
\toprule
\textbf{Categoria} & \textbf{Funcionalidade} & \textbf{Descrição Técnica} \\
\midrule
Cálculos & RLM Simples (Baseline) & Regressão Linear Múltipla padrão sem técnicas de ensemble \\
Cálculos & RLM+ Bagging & Ensemble com amostragem com reposição (bootstrap) \\
Cálculos & RLM+ Subagging+ CV & Ensemble sem reposição + otimização por validação cruzada 5-Fold \\
Configuração & Parâmetro $m$ (Bags) & Define número de amostras/ensembles para Bagging \\
Configuração & Parâmetro $k$ (Subagging) & Define fração ou quantidade para amostragem no Subagging \\
Visualização & Métricas de Desempenho & Exibição de $R^2$, RMSE, MAE no conjunto de teste \\
Visualização & Gráficos Comparativos & Scatter plots (previsto $\times$ real), resíduos e comparação entre modelos \\
Organização & Abas de Navegação & Separação lógica: ``Arquivos'', ``Configuração'', ``Resultados'', ``Gráficos'' \\
Exportação & Não identificada (legado) & Sistema original sem botões de download; corrigido na reengenharia \\
\bottomrule
\end{tabularx}
\end{table}

\subsection{Entradas e Saídas}
\begin{table}[H]
\centering
\caption{Entradas necessárias e saídas esperadas}
\begin{tabularx}{\textwidth}{l l X}
\toprule
\textbf{Tipo} & \textbf{Nome/Conteúdo} & \textbf{Descrição/Formato} \\
\midrule
\multicolumn{3}{c}{\textit{Entradas}} \\
Arquivo & \texttt{Xcal.txt} & Matriz de espectros de calibração (variáveis independentes) \\
Arquivo & \texttt{Ycal.txt} & Vetor de valores de referência para calibração \\
Arquivo & \texttt{Xteste.txt} & Matriz de espectros para predição/teste \\
Arquivo & \texttt{Yteste.txt} & Vetor de valores reais para validação \\
Parâmetro & $m$ e $k$ & Inteiro $\ge 1$ e Decimal $(0,1]$, respectivamente \\
\midrule
\multicolumn{3}{c}{\textit{Saídas}} \\
Numérico & $R^2$, RMSE, MAE & Tabela ou cards destacados \\
Gráfico & Previsto vs. Real & Scatterplot com reta de ajuste \\
Gráfico & Resíduos da Predição & Dispersão ou histograma \\
Comparativo & Desempenho entre modelos & Tabela ou gráfico de barras \\
Feedback & Mensagens de erro/validação & Alertas em destaque na interface \\
Status & Indicador de processamento & \texttt{st.progress} com etapas descritivas \\
\bottomrule
\end{tabularx}
\end{table}

\subsection{Classificação por Prioridade}
\begin{itemize}[leftmargin=1.5cm]
    \item \textbf{Essenciais (Must Have):} Upload e validação dos 4 arquivos; execução da RLM Simples; cálculo de métricas básicas; validação dimensional ($X$ vs $Y$); exibição clara dos resultados.
    \item \textbf{Desejáveis (Should Have):} Gráficos interativos (zoom, tooltip); exportação (CSV, TXT, PNG); tooltips explicativos para $m$ e $k$; pré-processamento básico; histórico de execuções na sessão.
    \item \textbf{Extras (Could Have):} Autenticação e projetos privados; integração com APIs externas; suporte a formatos (.csv, .xlsx, .mat); otimização automática de hiperparâmetros; dashboard comparativo multi-rodada.
\end{itemize}

\subsection{Requisitos Não Funcionais}
\begin{table}[H]
\centering
\caption{Especificação de requisitos não funcionais}
\begin{tabularx}{\textwidth}{l X X}
\toprule
\textbf{Categoria} & \textbf{Especificação} & \textbf{Critério de Aceite} \\
\midrule
Performance & Tempo de cálculo (RLM/Ensemble) & RLM $< 10\,$s; Ensemble $< 2\,$min (CV 5-Fold) \\
Performance & Renderização de gráficos & $< 3\,$s após conclusão do cálculo \\
Usabilidade & Acessibilidade (WCAG 2.1 AA) & Contraste, labels, navegação por teclado \\
Usabilidade & Clareza textual e recuperação de erro & Mensagens orientam correção, não apenas apontam falha \\
Segurança & Isolamento de sessão e sanitização & Dados não vazam; validação de extensão/conteúdo \\
Escalabilidade & Crescimento de usuários/dados & Container com load balancer; processamento em batch \\
\bottomrule
\end{tabularx}
\end{table}

\subsection{Fluxo Principal e Tratamento de Exceções}
O fluxo principal segue: \texttt{Upload $\rightarrow$ Validação dimensional $\rightarrow$ Configuração de $m$ e $k$ $\rightarrow$ Execução com CV 5-Fold $\rightarrow$ Exibição de resultados e gráficos}. Casos de exceção (arquivo inválido, parâmetros fora do intervalo, timeout) são tratados com validação \textit{fail-fast}, sliders com limites explícitos e mensagens orientativas, garantindo usabilidade e robustez.

\section{Projeto Lógico}

A reengenharia adotou o modelo clássico de desenvolvimento, promovendo a separação de responsabilidades e a eliminação de acoplamentos presentes no código legado.

\subsection{Arquitetura em Camadas}
\begin{table}[H]
\centering
\caption{Diagnóstico do código legado e arquitetura proposta}
\begin{tabularx}{\textwidth}{l X X}
\toprule
\textbf{Aspecto} & \textbf{Estado Atual (Legado)} & \textbf{Solução na Reengenharia} \\
\midrule
Arquitetura & Monolítica (UI + lógica + parsing) & Camadas separadas: Apresentação, Negócio, Dados \\
Validação & \texttt{try/except} genérico & Função \texttt{validate\_inputs()} com regras explícitas \\
Performance & Laços \texttt{for} em CV, sem cache & \texttt{@st.cache\_data}, vetorização NumPy, seed configurável \\
Tipagem/Doc & Ausente em funções críticas & Type hints + docstrings padrão NumPy \\
Reprodutibilidade & \texttt{seed=42} hardcoded & Parâmetro configurável + registro em logs \\
\bottomrule
\end{tabularx}
\end{table}

\subsection{Fluxo de Dados e Lógica de Negócio}
\begin{enumerate}[leftmargin=1.5cm]
    \item \textbf{Camada de Dados:} Parsing de arquivos \texttt{.txt} em \texttt{numpy.ndarray}, validação dimensional entre $X$ e $Y$, e sanitização de conteúdo.
    \item \textbf{Camada de Negócio:} Implementação vetorizada da RLM ($\hat{\beta} = X^+ Y$), Bagging/Subagging com reamostragem controlada, e validação cruzada 5-Fold para busca de hiperparâmetros ótimos ($m^*, k^*$).
    \item \textbf{Camada de Apresentação:} Interface Streamlit com abas, \texttt{st.session\_state} para persistência de resultados, e componentes Plotly para visualização interativa.
\end{enumerate}

\subsection{Estratégia de Otimização}
A validação cruzada 5-Fold particiona aleatoriamente o conjunto de calibração. Para cada par $(m, k)$, calcula-se o RMSE médio e desvio padrão nos folds. O par com menor RMSE-CV é selecionado e aplicado ao conjunto de teste independente, garantindo avaliação não enviesada.

\section{Projeto Físico}

\subsection{Estrutura de Diretórios e Responsabilidades}
\begin{table}[H]
\centering
\caption{Estrutura modular e responsabilidades}
\begin{tabularx}{\textwidth}{l X}
\toprule
\textbf{Módulo/Arquivo} & \textbf{Função Específica} \\
\midrule
\texttt{app.py} & Camada de apresentação (apenas UI + roteamento de abas) \\
\texttt{src/data\_loader.py} & Parsing de \texttt{.txt}, validação dimensional, dataclasses \\
\texttt{src/models.py} & Implementação de RLM, Bagging e Subagging (vetorizado) \\
\texttt{src/metrics.py} & Cálculo de RMSE, MAE e $R^2$ com tipos explícitos \\
\texttt{src/optimizers.py} & Validação cruzada 5-Fold e busca de hiperparâmetros \\
\texttt{tests/} & Suíte de testes unitários (\texttt{pytest}) para cada módulo \\
\bottomrule
\end{tabularx}
\end{table}

\subsection{Ambiente e Dependências}
\begin{itemize}[leftmargin=1.5cm]
    \item \textbf{Linguagem/ Framework:} Python 3.10+ e Streamlit 1.30+
    \item \textbf{Cálculo Numérico:} \texttt{numpy}, \texttt{pandas}
    \item \textbf{Visualização:} \texttt{plotly} (gráficos interativos)
    \item \textbf{Testes:} \texttt{pytest}, \texttt{pytest-cov}
    \item \textbf{Deploy:} Streamlit Community Cloud (gratuito) ou container Docker com \texttt{nginx} + \texttt{gunicorn} para escalabilidade.
\end{itemize}

\subsection{Persistência e Cache}
Utiliza-se \texttt{st.session\_state} para armazenar resultados entre interações, evitando reprocessamento. A renderização de gráficos é desacoplada do cálculo, e a validação dimensional ocorre antes da alocação em memória, prevenindo \texttt{MemoryError} em datasets espectrais grandes.

\section{Implementação}

A implementação materializa o projeto lógico e físico em código Python, mantendo fidelidade matemática ao original e ampliando usabilidade e reprodutibilidade.

\subsection{Destaques da Implementação (\texttt{app.py})}
\begin{lstlisting}[language=Python]
def perform_rlm(Xcal, Ycal, Xteste):
    """RLM simples via pseudoinversa."""
    beta = np.linalg.pinv(Xcal) @ Ycal
    return Xteste @ beta

def perform_bagging(Xcal, Ycal, Xteste, num_bags, with_replacement, fraction, seed=42):
    """RLM + Bagging / Subagging com vetorização e seed configurável."""
    n = Xcal.shape[0]
    sample_size = n if with_replacement else max(1, int(n * fraction))
    rng = np.random.default_rng(seed)
    Ypred = np.zeros(Xteste.shape[0])
    for _ in range(num_bags):
        idx = rng.integers(0, n, size=sample_size) if with_replacement else \
              rng.choice(n, size=sample_size, replace=False)
        beta = np.linalg.pinv(Xcal[idx]) @ Ycal[idx]
        Ypred += (Xteste @ beta) / num_bags
    return Ypred
\end{lstlisting}

\subsection{Interface e Fluxo de Execução}
\begin{itemize}[leftmargin=1.5cm]
    \item \textbf{Upload:} 4 arquivos \texttt{.txt} com validação imediata de formato e dimensões.
    \item \textbf{Configuração:} Inputs para $m$ e $k$ com parsing robusto e tooltips.
    \item \textbf{Execução:} \texttt{st.progress} com \textit{callback} atualizado a cada iteração de CV.
    \item \textbf{Resultados:} Cards estilizados, tabela \texttt{pandas} exportável para CSV, e log detalhado.
    \item \textbf{Gráficos:} Plotly com \texttt{Scatter}, linha de referência $y=x$, e tooltips interativos.
\end{itemize}

\subsection{Testes, Validação e Rastreabilidade}
\begin{table}[H]
\centering
\caption{Estratégia de testes e critérios de aceite}
\begin{tabularx}{\textwidth}{l X X}
\toprule
\textbf{Tipo de Teste} & \textbf{Escopo} & \textbf{Meta de Cobertura} \\
\midrule
Unitário & Funções isoladas (parse, metrics) & $\ge 80\%$ (\texttt{pytest-cov}) \\
Integração & Fluxo completo: Upload $\rightarrow$ CV $\rightarrow$ Gráficos & 100\% dos caminhos críticos \\
Performance & Tempo de resposta para datasets típicos & RLM $< 10\,$s; Ensemble $< 2\,$min \\
Aceite (UAT) & Validação por pesquisadores/alunos & SUS $\ge 75$ (Bom/Excelente) \\
\bottomrule
\end{tabularx}
\end{table}

\begin{table}[H]
\centering
\caption{Matriz de rastreabilidade de requisitos (RTM)}
\begin{tabularx}{\textwidth}{l l X}
\toprule
\textbf{ID Requisito} & \textbf{Fase Clássica} & \textbf{Artefato de Validação} \\
\midrule
RF01.1–01.5 & 1, 3, 5 & \texttt{data\_loader.py}, testes de upload \\
RF02.1–02.3 & 1, 2, 4 & \texttt{optimizers.py}, logs de execução \\
RF03.1–03.3 & 2, 3, 4 & \texttt{models.py}, testes unitários de métricas \\
RF04.1–04.3 & 2, 4, 5 & \texttt{metrics.py}, UAT visual (Plotly) \\
RF05.1–05.3 & 1, 3, 5 & UI/UX, feedback de erros, \texttt{st.progress} \\
\bottomrule
\end{tabularx}
\end{table}

\subsection{Comparativo Vaadin $\rightarrow$ Streamlit}
\begin{itemize}[leftmargin=1.5cm]
    \item \texttt{@Route("")} $\rightarrow$ \texttt{st.set\_page\_config()}
    \item \texttt{Tabs} $\rightarrow$ \texttt{st.tabs()}
    \item \texttt{Upload + MemoryBuffer} $\rightarrow$ \texttt{st.file\_uploader()}
    \item \texttt{CompletableFuture} $\rightarrow$ Execução síncrona com \texttt{st.progress}
    \item \texttt{RegressionService (Java)} $\rightarrow$ Funções Python com NumPy vetorizado
    \item \texttt{EJML SimpleMatrix} $\rightarrow$ \texttt{numpy.linalg.pinv()}
\end{itemize}

A implementação resultante atende aos requisitos funcionais e não funcionais levantados, reduz a curva de aprendizado para usuários não especialistas, e estabelece base técnica sólida para evoluções futuras (autenticação, pipelines automatizados, suporte a novos formatos).

\end{document}